\chapter{The detector description}
\label{ch:geometry}

\chaptergoals{3 of 11 in Part~II: the first package of the chain.}
{\cref{ch:bigpicture}.}
{build \file{ship\_geometry.db}, find any volume and its position, and name the volumes that will record hits.}

\section{GeoModel in a nutshell}

The SHiP detector is described with
\href{https://geomodel.web.cern.ch}{GeoModel}, the geometry toolkit developed
for ATLAS. A GeoModel description is ordinary C++ that builds a tree of
volumes:

\begin{itemize}
  \item a \emph{shape} (\texttt{GeoBox}, \texttt{GeoTubs}, \texttt{GeoTrap},
        \ldots) gives the extent;
  \item a \emph{logical volume} (\texttt{GeoLogVol}) adds a name and a
        material;
  \item a \emph{physical volume} (\texttt{GeoPhysVol}, or
        \texttt{GeoFullPhysVol} when its global position must be retrievable)
        places a logical volume, with a transform, inside its mother.
\end{itemize}

The tree is built once and written to a single SQLite file,
\file{ship\_geometry.db}. Every downstream package (simulation, event display,
reconstruction) reads that file, so they all see exactly the same detector.

In \repo{Geometry} each subsystem is a factory class in
\file{subsystems/<Name>/} whose \texttt{build()} returns the subsystem's
top volume; \texttt{SHiPGeometryBuilder} creates the world (the cavern) and
places each subsystem at its $z$ position. Materials are defined once, in
\texttt{SHiPMaterials}.

\section{Coordinates, units and names}
\label{sec:coordinates}

\begin{center}
\begin{tabularx}{0.9\linewidth}{lL}
\toprule
Origin   & front face of the target (the first tungsten disk) \\
$z$      & along the beam, positive downstream \\
$y$      & vertical, pointing up \\
$x$      & horizontal, completing a right-handed system \\
Units    & millimetres and radians (GeoModel's native units) \\
Beam axis& 1.7\,m above the cavern floor \\
\bottomrule
\end{tabularx}
\end{center}

Volume names follow a path convention: \texttt{/SHiP/<subsystem>/<part>}, for
example \texttt{/SHiP/trackers/straw\_gas} or
\texttt{/SHiP/timing\_detector/bar}. Downstream code selects volumes by
matching on these names, so they matter: you will use them in
\cref{ch:aegir} to tell Geant4 which volumes are sensitive.

\section{What is implemented}

\begin{center}
\small
\begin{tabularx}{\linewidth}{l l L}
\toprule
Subsystem & Status & Content \\
\midrule
Cavern            & complete          & world volume with the rock cavities subtracted \\
Target            & complete          & 33 tungsten disks in a grooved steel core, jacket and endcap \\
Muon shield       & approximate       & 6 stations; box approximations of the real shapes \\
Neutrino detector & approximate       & veto, Si/W target, HCAL with individual fibres \\
Magnet            & approximate       & iron yoke, box-shaped coils \\
Decay volume      & implemented       & frustum: SBT steel structure, liquid-scintillator (LAB) cells, helium \\
Upstream tagger   & simulation-ready  & 4.4\,$\times$\,6.4\,m plane of 47\,000 polystyrene tiles \\
Trackers          & complete          & 4 stations $\times$ 4 stereo views, 9600 straws \\
Timing detector   & complete          & 330 scintillator bars \\
Calorimeter       & simulation-ready  & ECAL and HCAL sampling layers \\
\bottomrule
\end{tabularx}
\end{center}

The status column is taken from the \repo{Geometry} README, where each
subsystem links to its own README with the volume tree and materials. The
nominal envelope of every subsystem (where it starts and ends in $z$, how wide
it is) is tabulated in \file{subsystem\_envelopes.csv} in the same repository.

\section{Building the geometry}

\begin{exercise}{Build \file{ship\_geometry.db}}
In your \file{Geometry} checkout:
\begin{shell}
pixi run geometry
ls -lh build/ship_geometry.db
\end{shell}
The \cmd{geometry} task compiles the package and runs
\cmd{build/apps/build\_geometry build/ship\_geometry.db}. Copy the result to
your working directory, where all later chapters expect it:
\begin{shell}
cp build/ship_geometry.db "$SHIP_TUTORIAL/work/"
\end{shell}
\end{exercise}

\begin{note}
The \texttt{shipgeometry} conda package installs \cmd{build\_geometry} too,
so in the \pkg{aegir} environment \cmd{pixi run build\_geometry} writes the
same file into the current directory without compiling anything. Building
from the \repo{Geometry} checkout is what you do once you change the
description.
\end{note}

\begin{exercise}{Validate it}
\begin{shell}
pixi run ./build/apps/validate_geometry build/ship_geometry.db
\end{shell}
The validator checks the database structure and prints the volumes, the
materials and the positions it finds. Skim the material list: which material
fills the decay volume, and which fills the straws?
\end{exercise}

\begin{solution}
The centre of the decay volume is helium; the straws contain the straw gas
volume \texttt{/SHiP/trackers/straw\_gas} inside a mylar wall
(\texttt{/SHiP/trackers/straw\_wall}). The SBT cells are LAB (liquid
scintillator).
\end{solution}

\section{Looking at the geometry}

\cmd{gmex}, the GeoModel explorer, opens a \file{.db} file directly:
\begin{shell}
pixi run gmex build/ship_geometry.db
\end{shell}
In the volume tree on the left, expand \texttt{/SHiP}, tick the subsystems
you want to see and use the mouse to rotate and zoom.

\begin{exercise}{Find the decay volume}
Open the geometry in \cmd{gmex}. Show only the decay volume, the trackers and
the calorimeter. Roughly where in $z$ does the decay volume start and end?
Then check your estimate against \file{subsystem\_envelopes.csv}.
\end{exercise}

\begin{solution}
From the envelope table: upstream tagger 32.52--32.92\,m, decay volume and
SBT 32.92--83.32\,m, tracker stations at 83.57--84.57, 85.57--86.57,
92.57--93.57 and 94.57--95.57\,m around the spectrometer magnet
(87.07--92.07\,m), timing detector 95.82--96.32\,m, ECAL and HCAL from
96.87 to 99.77\,m.
\end{solution}

The event display can print the same information as text, which is quicker
for scripting:

\begin{shell}
cd "$SHIP_TUTORIAL/sea-cucumber"
pixi run sc --geometry ../work/ship_geometry.db --inspect 4
pixi run sc --geometry ../work/ship_geometry.db --inspect 6 --inspect-match "*trackers*"
\end{shell}

The number after \cmd{--inspect} is the depth of the tree to walk;
\cmd{--inspect-match} restricts the listing to names matching a pattern.

\begin{exercise}{The sensitive volumes}
\label{ex:sensitive}
In its default scoring mode, Geant4 records positive-energy-deposit steps
only in volumes declared sensitive, and in
\cref{ch:aegir} you will declare them by name. Using \cmd{--inspect},
\cmd{gmex}, or the subsystem READMEs, find a name fragment that selects the
active material of each of these five detectors. Decide whether a single
substring can select it without also selecting passive material:
\begin{enumerate}[label=(\alph*)]
  \item the upstream background tagger (UBT) tiles;
  \item the surround background tagger (SBT) liquid-scintillator cells;
  \item the gas of the straw tubes;
  \item the calorimeter's scintillator layers;
  \item the timing-detector bars.
\end{enumerate}
\end{exercise}

\begin{solution}
Names as built by the factories in \repo{Geometry}:
\begin{enumerate}[label=(\alph*)]
  \item tiles are \texttt{/SHiP/upstream\_tagger/fine\_tile} and
        \texttt{/SHiP/upstream\_tagger/coarse\_tile}; the air envelopes
        holding them end in \texttt{\_env}.
  \item cells are named
        \texttt{/SHiP/decay\_volume/sbt/sensors\_\ldots\_C<n>}; the aluminium
        walls of the same containers end in \texttt{\_W<n>}.
  \item \texttt{/SHiP/trackers/straw\_gas}.
  \item scintillators are \texttt{/SHiP/calorimeter/wide\_pvt\_\{h,v\}\ldots}
        and \texttt{/SHiP/calorimeter/thin\_ps\_\{h,v\}\ldots}; absorbers are
        \texttt{lead\_plate} and \texttt{iron\_plate}.
  \item \texttt{/SHiP/timing\_detector/bar}.
\end{enumerate}
The point of the exercise is the discussion in \cref{sec:sensitive}: Geant4
matches \emph{substrings}, one fragment per detector, so for (a), (b) and (d)
there is no single fragment that selects only the active parts. The workflow
in \cref{ch:aegir} uses the subsystem prefixes and accepts the extra hits in
passive material.
\end{solution}

\begin{organisernote}
Exercise~\ref{ex:sensitive} has no clean answer for three detectors with the
current naming. Ask the Geometry and Shannon developers whether the factories
should add a common tag to the active volumes (e.g.\ \texttt{\_sens}) before
the tutorial; if they do, update \file{code/workflows/hnl\_sim.jsonnet} and
this solution.
\end{organisernote}
